\documentclass[11pt]{article}

\usepackage[margin=1in]{geometry}
\usepackage{amsmath, amssymb, amsthm}
\usepackage{graphicx}
\usepackage{booktabs}
\usepackage{hyperref}
\usepackage{float}
\usepackage{caption}
\usepackage{subcaption}
\usepackage{enumitem}
\usepackage[numbers]{natbib}

\title{Koopman Analysis of Chua's Circuit with EDMD}
\author{Nonlinear Dynamics and Chaos Project}
\date{}

% --- Convenience macros ---
\newcommand{\R}{\mathbb{R}}
\newcommand{\C}{\mathbb{C}}
\newcommand{\N}{\mathbb{N}}
\newcommand{\Koop}{\mathcal{K}}
\newcommand{\Hil}{\mathcal{F}}

\begin{document}

\maketitle

\begin{abstract}
The Koopman operator provides a way to study nonlinear dynamical systems through a linear, but generally infinite-dimensional, perspective \citep{koopman1931}. Instead of evolving the state directly, the Koopman framework evolves observables of the state. This is appealing in nonlinear dynamics because spectral ideas such as eigenvalues, eigenfunctions, and modes can then be used even when the original system is nonlinear. In this project we first introduce the Koopman operator and show a simple example in which a nonlinear system can be made closed by augmenting the state with one nonlinear observable. We then briefly describe Extended Dynamic Mode Decomposition (EDMD), which approximates the Koopman operator from data \citep{williams2015}. In the second half of the paper we apply EDMD to experimental data from Chua's circuit measured using an Arduino R4 Minima. We analyze four regimes --- fixed point, limit cycle, period-doubled oscillation, and double-scroll chaos --- and use Koopman eigenvalues, reconstructions, dictionary sweeps, and eigenfunction coloring to show what the Koopman viewpoint reveals about the route to chaos.
\end{abstract}

\section{Introduction}
One of the recurring themes in nonlinear dynamics is that even simple-looking differential equations can produce behavior that is difficult to understand directly in state space. Linear systems are much easier to analyze: eigenvalues describe growth, decay, and oscillation, and solutions can be decomposed into modal pieces. Nonlinear systems do not generally permit such a direct spectral analysis in their original coordinates.

The Koopman operator offers a different viewpoint \citep{koopman1931}. Rather than seeking a linear description of the state variables themselves, it studies how \emph{observables} evolve under the dynamics. Although the underlying system may be nonlinear, the Koopman operator acting on observables is linear. This move transfers some of the tools of linear algebra into nonlinear dynamics, provided one is willing to work in a larger space of functions.

This paper has two goals. First, we present the basic Koopman idea in a simple setting where the lifted system closes exactly. Second, we apply a data-driven approximation of Koopman theory to Chua's circuit, a canonical nonlinear electronic system that exhibits a progression from fixed-point behavior to periodic motion, period doubling, and chaos.

\section{Koopman Background}

\subsection{The Koopman Operator}
Consider a discrete-time dynamical system
\begin{equation}
    x_{n+1} = F(x_n),
\end{equation}
where $x_n \in \mathbb{R}^d$. Instead of following the state directly, we consider an observable $g : \mathbb{R}^d \to \mathbb{C}$. The Koopman operator $\mathcal{K}$ acts on observables by composition with the dynamics:
\begin{equation}
    (\mathcal{K}g)(x) = g(F(x)).
\end{equation}
This operator is linear, since for observables $g_1,g_2$ and scalars $a,b$,
\begin{equation}
    \mathcal{K}(a g_1 + b g_2) = a \mathcal{K}g_1 + b \mathcal{K}g_2.
\end{equation}

The key point is that the nonlinearity of $F$ has not disappeared; rather, it has been moved into the function space on which $\mathcal{K}$ acts. In general this space is infinite-dimensional. The benefit is that if we can identify useful Koopman eigenfunctions $\xi_k$, satisfying
\begin{equation}
    \mathcal{K}\xi_k = \lambda_k \xi_k,
\end{equation}
then each eigenfunction evolves by a simple scalar multiplication:
\begin{equation}
    \xi_k(x_{n+1}) = \lambda_k \xi_k(x_n).
\end{equation}
This gives a linear evolution law for quantities derived from the nonlinear state.

\subsection{A Simple Closed Koopman Representation}
To make the Koopman idea concrete, consider the continuous-time nonlinear system
\begin{equation}
    \dot{x}_1 = \mu x_1, \qquad \dot{x}_2 = \lambda (x_2 - x_1^2),
\end{equation}
for which the augmented observable vector
\begin{equation}
    \psi(x) = \begin{bmatrix} x_1 \\ x_2 \\ x_1^2 \end{bmatrix}
\end{equation}
obeys a closed linear system. Indeed,
\begin{align}
    \frac{d}{dt} x_1 &= \mu x_1, \\
    \frac{d}{dt} x_2 &= \lambda x_2 - \lambda x_1^2, \\
    \frac{d}{dt}(x_1^2) &= 2 x_1 \dot{x}_1 = 2\mu x_1^2.
\end{align}
Therefore
\begin{equation}
    \frac{d}{dt}
    \begin{bmatrix}
        x_1 \\ x_2 \\ x_1^2
    \end{bmatrix}
    =
    \begin{bmatrix}
        \mu & 0 & 0 \\
        0 & \lambda & -\lambda \\
        0 & 0 & 2\mu
    \end{bmatrix}
    \begin{bmatrix}
        x_1 \\ x_2 \\ x_1^2
    \end{bmatrix}.
\end{equation}
This is an exact finite-dimensional Koopman-invariant subspace. Although the original system is nonlinear in the state variables because of the $x_1^2$ term, the lifted observable vector evolves linearly. The example is useful because it shows what one hopes to achieve more generally: by choosing the right observables, nonlinear evolution in state space can become linear evolution in a lifted space.

This example also clarifies an important point. Koopman theory does not say that every nonlinear system has a low-dimensional exact linear representation. Rather, it says that nonlinear systems can be studied through linear evolution of observables, and in favorable cases a finite dictionary of observables may capture a useful approximation.

\subsection{Koopman Modes, Eigenvalues, and Interpretation}
If a vector of observables $\psi(x)$ can be expanded in terms of Koopman eigenfunctions, then the dynamics take the form of a modal decomposition. In discrete time, one may write schematically
\begin{equation}
    \psi(x_n) \approx \sum_{k} \lambda_k^n \, \xi_k(x_0) \, \phi_k,
\end{equation}
where $\lambda_k$ are Koopman eigenvalues, $\xi_k$ are eigenfunctions, and $\phi_k$ are mode vectors. The eigenvalues indicate decay, growth, or oscillation. Complex conjugate eigenvalue pairs correspond to oscillatory behavior, and for periodic dynamics one expects eigenvalues near the unit circle.

This spectral language is especially attractive in nonlinear dynamics because it lets us discuss fixed points, limit cycles, and more complicated attractors using a common linear vocabulary, even when the original equations are nonlinear.

\subsection{Brief Description of EDMD}
Since the Koopman Operator is infinite-dimensional, practical computation requires a
finite-dimensional approximation. EDMD~\cite{williams2015} works in a
\emph{dictionary} of $N_\psi$ observables
\begin{equation}
    \boldsymbol{\psi}(\mathbf{x}) =
    \begin{bmatrix} \psi_1(\mathbf{x}) & \psi_2(\mathbf{x}) & \cdots &
    \psi_{N_\psi}(\mathbf{x}) \end{bmatrix}^\top \in \C^{N_\psi}.
    \label{eq:dict}
\end{equation}
The lifted snapshot matrices are
\begin{equation}
    \Psi_X = \begin{bmatrix} \boldsymbol{\psi}(\mathbf{x}_1) & \cdots &
    \boldsymbol{\psi}(\mathbf{x}_M) \end{bmatrix} \in \C^{N_\psi \times M},
    \qquad
    \Psi_Y = \begin{bmatrix} \boldsymbol{\psi}(\mathbf{x}_2) & \cdots &
    \boldsymbol{\psi}(\mathbf{x}_{M+1}) \end{bmatrix} \in \C^{N_\psi \times M},
    \label{eq:snapshots}
\end{equation}
where $(\mathbf{x}_k, \mathbf{x}_{k+1})$ are consecutive snapshot pairs from
one or more trajectories.
EDMD then computes a matrix $K$ such that
\begin{equation}
    \Psi_Y \approx K \Psi_X.
\end{equation}
In our implementation this approximation is obtained by least squares with truncated SVD regularization. The eigenvalues and eigenvectors of $K$ approximate Koopman spectral data. The quality of the approximation depends heavily on the chosen dictionary: if the observables reflect the structure of the nonlinear system, the lifted linear model can be much more informative.

\section{Experimental System: Chua's Circuit}

\subsection{Why Chua's Circuit}
Chua's circuit is a standard example in nonlinear dynamics because it exhibits several qualitatively different behaviors as parameters are varied. In particular, it provides a physically accessible route from simple behavior to chaos. This makes it a natural test case for the Koopman viewpoint: we can ask how the spectral picture changes as the system passes through a fixed point, a stable oscillation, a period-doubled oscillation, and finally the chaotic double scroll.

At the center of the circuit is its nonlinear element, commonly called the Chua diode. Despite the name, this is not an ordinary semiconductor diode, but an active nonlinear resistor with a piecewise-linear current-voltage characteristic \citep{chua1986,matsumoto1985}. That nonlinearity is what gives the circuit its rich behavior. The slope changes in the Chua diode produce the switching structure that underlies the double-scroll attractor, and they also motivate the piecewise-linear observables used later in the EDMD dictionaries.

\subsection{Data Collection}
The experimental data in this project were collected from Chua's circuit using an Arduino R4 Minima. Chua's circuit is one of the canonical low-dimensional electronic systems exhibiting chaos and was developed as a simple electronic realization of rich nonlinear dynamics \citep{chua1986,matsumoto1985}. In this implementation, the inductor was realized as a gyrator built from two op amps, and the Chua diode was also implemented using two op amps to produce the required nonlinear piecewise-linear response. This allowed the full circuit to be constructed from active analog building blocks while preserving the characteristic dynamics of the standard Chua design.

To interface the analog circuit with the microcontroller, additional signal-conditioning stages were required. Summing amplifiers and diode clamps were used to shift and limit the measured voltages so that they lay within the readable range of the Arduino ADC. This step was necessary because the native circuit voltages are bipolar, while the Arduino input expects a bounded unipolar voltage range. After conditioning, the Arduino sampled three measured signals, denoted here by $x$, $y$, and $z$, corresponding to the circuit state coordinates used throughout the analysis. For each regime, $5000$ samples were recorded:
\begin{itemize}[leftmargin=1.5em]
    \item fixed point,
    \item limit cycle,
    \item period-doubled oscillation,
    \item double-scroll chaotic regime.
\end{itemize}
These measurements were saved as time series and then organized into snapshot pairs for EDMD.

\begin{figure}[H]
    \centering
    \fbox{\rule{0pt}{2.2in}\rule{0.9\textwidth}{0pt}}
    \caption{Placeholder for a figure of the Chua circuit and measurement setup.}
    \label{fig:circuit}
\end{figure}

\begin{figure}[H]
    \centering
    \fbox{\rule{0pt}{2.4in}\rule{0.9\textwidth}{0pt}}
    \caption{Placeholder for measured time series and/or phase portraits from the four experimental regimes.}
    \label{fig:measured-data}
\end{figure}

\section{EDMD Analysis of the Four Regimes}

\subsection{Choice of Dictionary}
Because EDMD depends strongly on the observable dictionary, part of the analysis focused on comparing different families of observables. The code includes linear, polynomial, radial basis function, and Chua-matched piecewise-linear observables. The Chua-specific observables are motivated by the piecewise-linear nonlinearity in the circuit, so they are expected to capture the dynamics more efficiently than a completely generic basis. This follows the broader EDMD philosophy that approximation quality depends not only on the data, but on whether the chosen observables are rich enough to represent the important Koopman-invariant structures \citep{williams2015}.

This comparison is important conceptually. Koopman theory is not a black-box replacement for nonlinear dynamics. Rather, it is a framework in which physical insight enters through the observables one chooses to evolve.

\subsection{Dictionary Sweep: Accuracy Versus Cost}
To quantify the effect of dictionary choice, we performed a sweep over dictionary families and sizes and measured the relative one-step error in lifted space together with computation time. The expected trend is that richer dictionaries often reduce approximation error, but they also increase computational cost and can introduce conditioning issues. This tradeoff is visible in the implementation and is an important practical lesson of EDMD.

\begin{figure}[H]
    \centering
    \fbox{\rule{0pt}{2.5in}\rule{0.9\textwidth}{0pt}}
    \caption{Placeholder for dictionary sweep results showing how improved observable sets can reduce relative error while increasing computation time.}
    \label{fig:dictionary-sweep}
\end{figure}

\subsection{Koopman Eigenvalues Across Regimes}
For each regime, EDMD produces a finite-dimensional Koopman matrix whose eigenvalues approximate the spectral behavior of the system. These eigenvalues are especially informative when compared across regimes.

For a fixed point, one expects non-oscillatory behavior to dominate, so the spectrum should be concentrated near the real axis. For a stable limit cycle, complex conjugate pairs are expected because oscillation in real-valued systems is represented by conjugate modes. If the oscillation persists without decay, the corresponding eigenvalues should lie near the unit circle in discrete time. For a period-doubled orbit, one expects new spectral content associated with the subharmonic frequency, so additional conjugate pairs should appear at lower frequencies. In the chaotic double-scroll regime, the spectrum is generally more complicated, but one still hopes to identify coherent oscillatory structure through leading eigenvalues and eigenfunctions.

\begin{figure}[H]
    \centering
    \fbox{\rule{0pt}{2.5in}\rule{0.9\textwidth}{0pt}}
    \caption{Placeholder for the Koopman eigenvalues computed in each regime. Complex conjugate pairs are expected for oscillatory dynamics, especially in the limit-cycle and period-doubled regimes.}
    \label{fig:eigenvalues}
\end{figure}

\subsection{Period-Doubled Reconstruction from Koopman Modes}
One of the strongest demonstrations in this project is the Koopman reconstruction of the period-doubled regime. The analysis script identifies bounded modes and reconstructs the measured state from selected Koopman components. In the period-doubled case, the script deliberately switches to a richer polynomial dictionary for reconstruction because the smaller piecewise-linear dictionary is not expressive enough to capture the half-frequency eigenfunction associated with the doubled orbit.

This is a particularly instructive example because it shows how the spectral picture reflects the nonlinear bifurcation. The dominant oscillatory content is no longer just the original fundamental frequency. Instead, the reconstruction reveals the importance of the subharmonic mode, corresponding to the fact that the system now repeats only after two passes rather than one. In the code, the selected modes are also projected onto the unit circle during reconstruction to remove artificial numerical decay of sustained oscillations. This makes the physical interpretation of the periodic orbit clearer.

\begin{figure}[H]
    \centering
    \fbox{\rule{0pt}{2.8in}\rule{0.9\textwidth}{0pt}}
    \caption{Placeholder for the Koopman mode reconstruction of the period-doubled regime, together with an explanation of the mode selection and unit-circle projection implemented in \texttt{run\_chua.m}.}
    \label{fig:period-doubled-reconstruction}
\end{figure}

\subsection{Double Scroll and Eigenfunction Coloring}
The chaotic double-scroll regime is the most visually striking part of Chua's circuit, but it is also where the Koopman viewpoint becomes especially insightful. In raw state space the trajectory wanders unpredictably between two lobes, and no individual voltage---$x$, $y$, or $z$---cleanly distinguishes them. The two lobes overlap in every linear projection: a threshold on $x$ alone, for example, will incorrectly assign points near the center of the attractor because both lobes pass through that region. This is a fundamental limitation of working directly with state coordinates in a chaotic system.

The Koopman framework offers a different approach. Rather than looking for a separator in the original state space, we look for an eigenfunction of the Koopman operator, a nonlinear observable that evolves by simple multiplication at each timestep. If $\xi_k$ denotes such an eigenfunction, it satisfies
\begin{equation}
    \xi_k(x_{n+1}) = \lambda_k \,\xi_k(x_n),
    \label{eq:eigenfunction}
\end{equation}
so that $\xi_k(x_n) = \lambda_k^n \,\xi_k(x_0)$. The state $x_n$ itself evolves chaotically, but the scalar $\xi_k(x_n)$ executes a rigid rotation in the complex plane at the angular frequency $\arg(\lambda_k)$ per timestep. This is precisely the Koopman linearity guarantee: nonlinear dynamics in state space become linear dynamics in observable space.

\paragraph{Computing eigenfunctions from data.}
EDMD produces the finite Koopman matrix $K$ and its eigenvector matrix $\Phi$ (columns are right eigenvectors). The corresponding left eigenvectors, collected as rows of $W = \Phi^{+}$, provide the dual basis. The $k$-th Koopman eigenfunction evaluated along the trajectory is then
\begin{equation}
    \xi_k(x_n) = w_k^\top \psi(x_n),
\end{equation}
where $w_k$ is the $k$-th row of $W$ and $\psi(x_n)$ is the lifted dictionary vector at snapshot $n$. Evaluating this for every snapshot gives the full time history of $\xi_k$ along the measured attractor. The eigenfunction selected for visualization is the bounded oscillatory mode with the highest variance in $\mathrm{Re}(\xi_k)$, which is the mode that most strongly discriminates different regions of the attractor. For this regime a degree-3 polynomial dictionary augmented with the four piecewise-linear Chua kink terms was used, since a small dictionary does not provide enough expressive power to resolve the eigenfunction structure of the double scroll.

\paragraph{Interpreting the four panels.}
The visualization in Figure~\ref{fig:double-scroll-coloring} presents four complementary views of the same eigenfunction.

The top two panels color every attractor point by $\arg(\xi_k)$ using a cyclic (HSV) colormap, shown in the $x$-$y$ and $x$-$z$ projections respectively. Because $\xi_k$ advances by a fixed angle each step, points visited at similar phases of the rotation receive similar colors. If the eigenfunction genuinely reflects the attractor's geometry, the color should sweep smoothly around each lobe rather than appearing randomly speckled. Smooth, coherent color bands confirm that $\xi_k$ is acting as an intrinsic phase coordinate on the attractor---not a noise artifact---and that the chosen mode captures persistent oscillatory structure even within the chaos.

The bottom-left panel reduces the complex eigenfunction to a binary label: points where $\mathrm{Re}(\xi_k) \geq 0$ are shown in blue, and those where $\mathrm{Re}(\xi_k) < 0$ in red. This is the most direct test of the eigenfunction's utility as a lobe separator. If it succeeds, blue and red should cleanly partition the two scrolls in the full three-dimensional state space, achieving in a single learned observable what no voltage threshold can accomplish. The key is that $\xi_k$ is a nonlinear function of all three state variables: the EDMD procedure has found, from data alone, the combination of polynomial and piecewise-linear terms that most cleanly partitions the attractor.

The bottom-right panel shows $\mathrm{Re}(\xi_k)$ and $\mathrm{Im}(\xi_k)$ as time series. Despite the irregular switching visible in the raw voltages, the eigenfunction trace should show approximately sinusoidal oscillation whose real part changes sign each time the trajectory transitions between lobes. This is the time-domain signature of the linear evolution in~\eqref{eq:eigenfunction}: the chaotic switching in physical coordinates maps to a regular phase rotation in the eigenfunction coordinate.

\paragraph{Why this demonstrates the advantage of Koopman.}
The central challenge for any analysis method applied to chaos is that the dynamics are sensitive to initial conditions and superficially irregular. State-space methods that rely on linear combinations of voltages are limited because nonlinear geometric features---such as lobe membership---have no simple linear signature. The Koopman approach circumvents this by explicitly searching the space of nonlinear observables for one that evolves linearly. The eigenfunction coloring demonstrates this advantage concretely: it transforms the question ``which lobe is the trajectory in?'' from an intractable problem in state space into a trivial threshold on a single learned coordinate. Moreover, the eigenvalue $\lambda_k$ associated with that coordinate tells us the frequency at which the trajectory winds around the attractor, providing quantitative dynamical information that is completely invisible in the raw time series.

\begin{figure}[H]
    \centering
    \fbox{\rule{0pt}{2.8in}\rule{0.9\textwidth}{0pt}}
    \caption{Double-scroll attractor colored by Koopman eigenfunction $\xi_k$ (mode selected by oscillatory variance). \emph{Top row:} attractor points colored by $\arg(\xi_k)$ in the $x$-$y$ and $x$-$z$ projections; coherent color bands confirm that $\xi_k$ is a genuine phase coordinate on the attractor. \emph{Bottom left:} three-dimensional attractor split by $\mathrm{sign}(\mathrm{Re}(\xi_k))$; blue and red cleanly separate the two scrolls, a partition that no single voltage threshold achieves. \emph{Bottom right:} time trace of $\mathrm{Re}(\xi_k)$ (blue) and $\mathrm{Im}(\xi_k)$ (red), showing near-sinusoidal evolution of the eigenfunction despite irregular chaotic switching in the raw state.}
    \label{fig:double-scroll-coloring}
\end{figure}

\section{Discussion}
Several points emerge from this project. First, the Koopman framework provides a natural spectral language for nonlinear dynamics. Even though the systems under study are nonlinear, their behavior can be described through eigenvalues, eigenfunctions, and modes once appropriate observables are chosen.

Second, dictionary choice is central. The dictionary sweep shows that better observables can substantially improve the EDMD approximation, but not for free: richer dictionaries increase computation time and can worsen numerical conditioning. This is both a practical issue and a conceptual one, since it emphasizes that Koopman analysis depends on choosing observables that reflect the underlying physics.

Third, the Chua circuit data illustrate a meaningful progression through the route to chaos. The fixed-point regime is spectrally simple, the limit cycle exhibits conjugate-pair structure associated with sustained oscillation, the period-doubled regime reveals subharmonic content, and the double-scroll regime shows that Koopman eigenfunctions can still uncover coherent structure even in chaos.

\section{Conclusion}
This project used the Koopman operator as a bridge between linear spectral ideas and nonlinear dynamical behavior. A simple background example showed how augmenting the observable space can produce a closed linear description when the right nonlinear term is included. EDMD then provided a practical way to approximate Koopman spectral quantities from experimental data \citep{koopman1931,williams2015}.

Applied to Chua's circuit, this framework gave a unified way to study four regimes measured with an Arduino R4 Minima: fixed point, limit cycle, period doubling, and double-scroll chaos. The most compelling result is that the Koopman viewpoint does not merely fit the data; it organizes the behavior. In the periodic regimes it yields interpretable oscillatory modes, and in the chaotic regime it produces eigenfunctions that reveal hidden coordinates on the attractor. This makes the Koopman approach especially valuable for a nonlinear dynamics setting, where the central challenge is often not only to simulate nonlinear motion, but to understand its structure.

\bibliographystyle{plainnat}
\bibliography{bib}

\end{document}
